\chapter{总结与展望}

\section{全文工作总结}
本文围绕流程工业场景下的领域自适应故障诊断问题，面向 Tennessee Eastman Process Domain Adaptation 数据集，完成了从研究现状梳理、统一实验基准构建、典型基线比较到 RCTA 方法拆解设计与递进验证规划的一系列工作。全文首先在绪论中梳理国内外研究现状，并将研究挑战与后续创新章节建立对应关系；随后说明跨工况迁移诊断的统一任务边界、共享骨干网络、基础训练目标和典型领域自适应基线；在此基础上，将 RCTA 拆分为教师--学生时序一致性、可靠性门控伪标签、双原型记忆类条件约束和课程化优化四个创新环节；最后在实验章节中按照 A、A+B、A+B+C、A+B+C+D 的顺序组织递进验证框架。

相较于原有章节组织，当前骨架将文献综述前移至绪论，将原问题定义、数据集、实验设定和基线结果内化到基准章节与实验章节中，并把 RCTA 的关键模块拆分为独立章节。这样的结构更有利于体现“现有痛点--方法创新--实验验证”之间的对应关系。

\section{主要结论}
在当前骨架版本中，本文可确认的阶段性结论主要包括以下四点。

\begin{enumerate}
\item TEP 多工况故障诊断可以被清晰表述为单源到单目标的无监督领域自适应问题，并能够在统一数据划分、统一骨干网络和统一模型选择准则下开展可复现比较。
\item 典型领域自适应基线方法能够为 RCTA 的设计提供参照，其中全局分布对齐、对抗式对齐和最优传输方法各有优势，但在多类别工业时序任务中仍需要进一步处理目标域伪标签噪声和类别结构混叠问题。
\item RCTA 的四个创新环节分别回应四类痛点：教师--学生时序一致性用于稳定目标域预测来源，可靠性门控用于筛选可信伪标签，双原型记忆用于显式塑造类条件结构，课程化优化用于控制各辅助目标的进入时机。
\item 倒数第二章已经按照创新点顺序预留了递进实验框架，后续可直接回填 A、A+B、A+B+C、A+B+C+D 的定量结果、混淆矩阵和 t-SNE 可视化，从而形成完整的验证闭环。
\end{enumerate}

\section{不足与未来展望}
尽管本文已经在单源到单目标跨工况诊断场景上形成了较完整的研究链路，但仍存在若干值得进一步完善之处。

\begin{enumerate}
\item 当前正文聚焦四个关键单源场景，尚未对更大范围的迁移对进行系统统计，因此对于不同工况组合下方法收益边界的认识仍可继续拓展。
\item RCTA 的完整验证结论仍待最终结果回填，后续需要进一步补充严格的模块消融、参数敏感性分析和统计显著性检验，以明确各模块的独立贡献与相互作用。
\item 当前分析主要集中在离线无监督领域自适应设置，对更复杂迁移场景、在线适应、持续诊断以及工程部署约束等问题尚未展开深入研究。
\end{enumerate}

围绕上述不足，后续工作可以从以下几个方向继续推进：其一，扩展到更多迁移方向和更复杂场景，系统考察 RCTA 在更广泛工业诊断任务中的稳定性；其二，围绕可靠性门控、双原型记忆、一致性正则和课程化调度开展严格消融实验，进一步厘清方法机制；其三，引入更丰富的可解释性分析，如类别级误差分解、原型演化轨迹和伪标签质量统计，以增强方法分析深度；其四，探索与工业先验知识、在线更新机制和轻量化部署约束结合的混合方法，为后续工程应用奠定基础。

\section{本章小结}
总体来看，本文已经围绕 TEP 跨工况故障诊断任务形成了一条更清晰的论文主线：绪论中交代研究现状与挑战，基准章节统一实验前提，方法章节按创新点逐一展开，实验章节按同一顺序递进验证，最后总结全文并展望后续工作。当前版本重点完成骨架调整，后续可在既有结构中继续补入最终实验结果与精修分析。
